%%%%%%%%%%  scan idx 690  |  folio 677  %%%%%%%%%%
\SGAPrintPage{677}%
of $Y$ in $X$, assumed everywhere to have rank $d$. In loc. cit., this formula follows from a very simple computation of $\underline{\Tor}$ sheaves. By what we saw in 4.2, the preceding formula implies the same formula in $A(Y)$, but only \uline{modulo torsion}. It is quite extraordinary that formula (4.6) in $A(Y)$ itself has not been proved at present\SGAAuthorFootnote{*}{cf. the note on page 23.\newline (**)}, nor even the special case obtained by taking $y=1$, which gives the self-intersection of $Y$:
\sgaStableTarget{sga6.e14.eq.4-7}\begin{equation}\tag{4.7}
 f^*(\cl(Y))=c_d(\underline{N}),
\end{equation}
even when $k$ is the field of complex numbers!

An analogous problem arises when one attempts to compute the Chow ring of the variety $X'$ obtained by blowing up $X$ along $Y$, imitating the method used to compute $\operatorname{Gr}^\bullet(X')_\Q$ [Exp. VII 4.8], which also works (without tensoring with $\Q$) for the computation of $\Gtop^\bullet(X')$. The key formula for this determination (with the notation of loc. cit.) is
\sgaStableTarget{sga6.e14.eq.4-8}\begin{equation}\tag{4.8}
 f^*(i_*(y))=j_*\bigl(g^*(y)c_{d-1}(\check{\underline F})\bigr),
\end{equation}
a formula currently established in $\Gtop^\bullet(X')$ (for $y\in\Gtop^\bullet(Y)$) by a computation of $\underline{\Tor}$ (Exp. VII 3.4), whose validity in $A(X')$ (for $y\in A(Y)$) remains to be established.

Observe that even the $\ell$-adic analogues of formulas (4.6), (4.7), (4.8) have not been proved at present; when $y$ is the cohomology class of an algebraic cycle, (4.6) and (4.8) hold in any case mod. torsion.

\sgaStableTarget{sga6.e14.para.4-5}\textbf{4.5.} In this section and the following ones, we shall give examples in which, in diagram (4.1), the homomorphism $j$, respectively the homomorphism $\phi$, is not an isomorphism. For a given $X$, it clearly amounts to the same thing to say that
\[
 j:\operatorname{Gr}^\bullet(X)\longrightarrow \Gtop^\bullet(X)
\]
is an isomorphism, or a monomorphism, or an epimorphism. Moreover, the image of $j$ is the subring of $\Gtop^\bullet(X)$ generated by the Chern classes of locally free sheaves on $X$. When $K^\bullet(X)$ is generated as an algebra over $\Z$ by the classes of invertible sheaves on $X$, one sees immediately that the subring of $\Gtop^\bullet(X)$ under consideration is also the subring generated by $\Gtop^1(X)$, the component of degree $1$ of $\Gtop^\bullet(X)$.

%%%%%%%%%%  scan idx 691  |  folio 678  %%%%%%%%%%
\SGAPrintPage{678}%
Now take for $X$ a \uline{flag variety}
\[
 X=G/B,
\]
where $G$ is a simply connected semisimple algebraic group over the algebraically closed field $k$, and $B$ is a Borel subgroup [5]. When $k=\C$, the Bruhat cell decomposition of $X$ [5, Exp. 13] implies that the canonical homomorphism
\sgaStableTarget{sga6.e14.eq.unnamed}\[
 (*)\qquad K^\bullet(X)\longrightarrow K(X(\C))
\]
is an isomorphism, as is moreover the canonical homomorphism
\sgaStableTarget{sga6.e14.eq.unnamed.occ2}\[
 (**)\qquad A(X)\longrightarrow H^*(X(\C),\Z),
\]
where $X(\C)$ is the compact space of complex points of $X$, endowed with its usual topology induced by the topology of $\C$. On the other hand, L. HODGKIN [12] has proved that $K(X(\C))$ is indeed generated by the classes of invertible sheaves on $X$. The same therefore holds for $K^\bullet(X)$ itself, and it is moreover not difficult to deduce, by a specialization technique using split semisimple groups over $\Z$ (SGA 3), that the same result holds over every algebraically closed field $k$. Moreover, the cell decomposition of $X$ (or the bijectivity of $(**)$ when $k=\C$) readily implies that $A(X)$ is torsion-free, so by 4.2 the homomorphism $\phi:A(X)\longrightarrow\Gtop^\bullet(X)$ is an isomorphism. These facts show that $j$ is an isomorphism if and only if $A(X)$ is generated by its elements of degree $1$, or again, when $k=\C$, if and only if the integral cohomology of $X(\C)$ is generated by the elements of degree $2$, i.e. if the classifying space $B_{G(\C)}$ is torsion-free. It is also known [15, p. 217 th. B (1)$\Leftrightarrow$(2), and 2.5] that this is so if and only if $G$ is "torsion-free", or equivalently if $G$ is isomorphic to a product of groups $\operatorname{S\ell}(n)$ and $\operatorname{Sp}(m)$. Thus it suffices to take for $G$ a group $\operatorname{SO}(n)$ with $n\geq 7$, or an exceptional simple group, to obtain an example in which $j$ is not an isomorphism.

\sgaStableTarget{sga6.e14.para.4-6}\textbf{4.6.} Observe that, in this example, there is no doubt that the "right" ring is the torsion-free ring $A(X)\simeq\Gtop^\bullet(X)$, with the cells of the Bruhat decomposition as a basis, a ring also isomorphic to the integral cohomology of $X(\C)$ when $k=\C$; on the other hand, the ring $\operatorname{Gr}^\bullet(X)$ appears to have torsion of a rather pathological nature. This example therefore confirms the view that the invariant $\operatorname{Gr}^\bullet(X)$ is of little
%%%%%%%%%%  scan idx 692  |  folio 679  %%%%%%%%%%
\SGAPrintPage{679}%
interest except modulo torsion, i.e. after tensoring with $\Q$. At the same time, this example shows that only after tensoring with $\Q$ is $\operatorname{Gr}^\bullet(X)$ covariant with respect to proper morphisms, even when restricted to closed immersions of smooth quasi-projective schemes. Indeed, for $X$ such that $j$ is not an isomorphism, there is in $X$ an integral closed subscheme $Y$ such that $\mathrm{c}\ell^\bullet(Y)\in K^\bullet(X)$ does not lie in $\operatorname{Filt}^i(X)$, where $i=\operatorname{codim}(Y,X)$. We may take $Y$ of minimal dimension, so that $\operatorname{Filt}^{i+1}(X)=\operatorname{Filt}_{\mathrm{top}}^{i+1}(X)$. With $Z$ the singular locus of $Y$ and $X'=X-Z$, $Y'=Y-Z$, we conclude that $\mathrm{c}\ell^\bullet(Y')\in K^\bullet(X')$ does not lie in $\operatorname{Filt}^i(X')$ (since $K^\bullet(X)\to K^\bullet(X')$ is surjective and compatible with the augmentations, and its kernel is $\operatorname{Im}\bigl(K(Z)\longrightarrow K(X)\bigr)\subset\operatorname{Filt}_{\mathrm{top}}^{i+1}(X)=\operatorname{Filt}^{i+1}(X)$). Now consider the inclusion
\[
 f:Y'\longrightarrow X',
\]
which is a closed immersion of codimension $i$ of smooth quasi-projective schemes over $k$, \uline{such that}
\[
 f_*:K^\bullet(Y')\longrightarrow K^\bullet(X')
\]
\uline{is not compatible with the} $\lambda$-\uline{filtrations} (\uline{up to a shift by} $i$), since $f_*(1)$ does not lie in $\operatorname{Filt}^i(X')$.

\sgaStableTarget{sga6.e14.para.4-7}\textbf{4.7.} Keeping the notation of 4.5, observe now that (Exp. IX) shows that the cited theorem of HODGKIN is equivalent to the existence of an isomorphism
\[
 K^\bullet(G)\simeq\Z,\qquad\text{therefore}\qquad \Gtop^\bullet(G)\simeq\Z,
\]
the homomorphism $K^\bullet(G/B)\to K^\bullet(G)$ being a priori surjective, with kernel the ideal generated by the $\mathrm{c}\ell^\bullet(\underline{L})-1$, where $\underline{L}$ ranges over the invertible Modules on $X=G/B$ associated with the representations $B\to\underline{G}_m$ (which in fact give all the invertible modules on $X$ [5, Exp. 15]). On the other hand, an analogous argument shows that
\[
 A(G/B)\longrightarrow A(G)
\]
is also surjective and that its kernel is the ideal generated by the elements of $A^1(G/B)$ defined by the representations $B\to\underline{G}_m$, thus in fact by all of $A^1(G/B)$.\SGAAuthorFootnote{*}{Cf. IX 3.4, 2.2.} This shows that $A(G)$ reduces to $A^0(X)=\Z$ if and
%%%%%%%%%%  scan idx 693  |  folio 680  %%%%%%%%%%
\SGAPrintPage{680}%
only if "$G$ is torsion-free", i.e. $G$ is isomorphic to a product of groups $\operatorname{S\ell}(n)$ and $\operatorname{Sp}(m)$. Thus in this case, and in this case alone, the homomorphism
\[
 \phi:A(G)\longrightarrow\Gtop^\bullet(G)
\]
is an isomorphism. Thus, when $G$ has torsion, the underlying scheme provides an example of a nonbijective $\phi$.

\sgaStableTarget{sga6.e14.para.4-8}\textbf{4.8.} Similarly, one finds an example in which the homomorphism $\gamma$ of (4.1)
\sgaStableTarget{sga6.e14.eq.4-9}\begin{equation}\tag{4.9}
 \gamma:A(G)\longrightarrow H^{2*}(G,\underline{Z}_\ell(*))
\end{equation}
does not factor through $\Gtop^\bullet(G)$, or, equivalently here, an example in which this homomorphism is nonzero in degrees $>0$. Restricting ourselves (as suitable specialization theorems allow) to the case $k=\C$, we are reduced to studying the homomorphism
\sgaStableTarget{sga6.e14.eq.4-10}\begin{equation}\tag{4.10}
 A(G)\longrightarrow H^*(G(\C),\Z),
\end{equation}
which is moreover deduced by passage to the quotient from the canonical homomorphism
\[
 A(G/B)\simeq H^*(G/B(\C),\Z)\longrightarrow H^*(G(\C),\Z).
\]
It is tempting to conjecture that the homomorphism (4.10) is in fact injective, since the torsion primes of $G(\C)$ are known to be exactly those dividing the order of the finite group $A^+(G)=\coprod_{i>0}A^i(G)$. Were this so, as soon as $G$ had torsion the homomorphism (4.10) would not factor through $\phi$, and thus, for suitable $\ell$ (more precisely, for $\ell$ dividing the order of $A^+(G)$, i.e. $\ell$ a torsion coefficient of $G$), the same would hold for (4.9). One can at least assert that, for the group $G=\underline{G}_2$ over the field $\C$, the homomorphism (4.10) is nonzero in degree $>0$; more precisely, there is an element of $A^3(G)$ whose image in $H^6(G(\C),\Z)$ has order $2$; consequently, the homomorphism (4.9) does not factor through $\phi$ when $\ell=2$. The fact used concerning $\underline{G}_2$, which amounts to the existence of an element of $H^6(G/B(\C),\Z)$ whose image in $H^6(G(\C),\Z)$ is nonzero, is found in [3], using the familiar dictionary between the languages of complex semisimple algebraic groups and compact Lie groups; it was kindly communicated to me by J.P. SERRE.

%%%%%%%%%%  scan idx 694  |  folio 681  %%%%%%%%%%
\SGAPrintPage{681}%
\sgaStableTarget{sga6.e14.sec.5}\section*{5. Relations between \(\operatorname{Gr}^\bullet(X)\) and \(H^{2*}(X,\underline{Z}_\ell(*))\)}

\sgaStableTarget{sga6.e14.para.5-1}\textbf{5.1.} The example given in 4.8 does not rule out the possibility of finding, for an arbitrary scheme $X$, a ring homomorphism
\[
 g^\bullet:\operatorname{Gr}^\bullet(X)\longrightarrow H^{2*}(X,\underline{Z}_\ell(*)),
\]
such that, for every locally free Module $\underline{E}$ on $X$, $g^i$ takes the Chern class $c_i(\underline{E})$ of Exp. V to the cohomological Chern class of SGA 5 VII. Of course, $\ell$ here denotes a prime distinct from the residue characteristics of $X$. Since the ring $\operatorname{Gr}^\bullet(X)$ is generated by the $c^i(\underline{E})$, one sees in any case that $g^\bullet$ is uniquely determined by the requirement of compatibility with Chern classes.

\sgaStableTarget{sga6.e14.para.5-2}\textbf{5.2.} Using the Chern homomorphism
\[
 K^\bullet(X)\longrightarrow\Ch\bigl(H^{2*}(X,\Zell(*))\bigr),
\]
and passing to the associated graded groups for the filtration of the first member and the degree filtration of the second, respectively, one obtains a homomorphism of graded groups (but not of rings!)
\[
 \psi:\operatorname{Gr}^\bullet(X)\longrightarrow H^{2*}(X,\Zell(*)),
\]
satisfying the relation
\[
 \psi(c_i(\underline{E}))=(-1)^{i-1}(i-1)!\,c_i^\ell(\underline{E}),
\]
where $c_i^\ell(\underline{E})$ denotes the $\ell$-adic Chern class. By division starting from the $\psi^i$, one obtains a ring homomorphism:
\[
 g_\Q:\operatorname{Gr}^\bullet(X)\longrightarrow
 H^{2*}(X,\Qell(*))=
 H^{2*}(X,\Zell(*))\otimes_{Z_\ell}\Qell,
\]
taking Chern classes to Chern classes, such that the image of the first member is contained in that of $H^{2*}(X,\Zell(*))$; consequently, the question of the existence of $g^\bullet$ has an affirmative answer upon replacing $H^{2*}(X,\Zell(*))$ by its quotient by its torsion subgroup. The first doubtful case in defining the components $g^i$ of the hypothetical homomorphism $g$ is that of $g^3$; a priori, one succeeds only in defining $2g^3=\psi^3$. The answer to the existence question for $g^3$ would be negative if, for example, one could find three invertible sheaves $\underline{L}_1,\underline{L}_2,\underline{L}_3$ on $X$ satisfying the relation
\[
 (\underline{L}_1-1)(\underline{L}_2-1)(\underline{L}_3-1)=0\quad\text{in }K^\bullet(X),
\]
